\chapter{Lösungsidee}

In die \emph{Lösungsidee} gehören, durch Absätze getrennt:

\begin{itemize}
 \item Beschreibung der Vorgehensweise, die für die Lösung des Problems verwendet werden soll
 \item Erläuterung der wichtigsten Algorithmen und Datenstrukturen, die verwendet werden
 \item Vorstellung alternativer Herangehensweisen/Algorithmen/Datenstrukturen mit kurzer Diskussion von deren Vor- und Nachteilen
\end{itemize}

\begin{figure}[ht]
 \centering
 \includegraphics[scale=0.25]{bilder/LogoCZG.pdf}
 \caption{Das Logo unserer Schule}
 \label{abb:beispiel}
\end{figure}

In der gesamten Arbeit soll auf Abbildungen verwiesen werden (siehe Abbildung \ref{abb:beispiel}) und gegebenenfalls hilfreiche Zitate eingefügt werden.\newline
Wörtliche Rede wird in Latex zwischen \verb|"`| und \verb|"'| gesetzt. 

\begin{quote}
 "`\blindtext"'\footnote{Heinemann, 1998, S.\,280}
\end{quote}
